% !TeX program = latexmk
% !TeX options = -xelatex -synctex=1 -interaction=nonstopmode -halt-on-error -file-line-error "%DOC%"
% !TeX encoding = UTF-8
% 编译：latexmk -xelatex -synctex=1 -interaction=nonstopmode -halt-on-error 第二问模型.tex
% 交叉引用依赖 .aux 文件，需多轮编译；latexmk 会自动运行至引用稳定。
\documentclass[UTF8,a4paper,zihao=-4,fontset=none]{ctexart}
\usepackage[a4paper,margin=2.45cm,headheight=15pt,footskip=1.15cm]{geometry}
\setCJKmainfont{FandolSong-Regular.otf}[cmap=UniGB-UTF16-H,BoldFont=FandolHei-Regular.otf,ItalicFont=FandolKai-Regular.otf]
\setCJKsansfont{FandolHei-Regular.otf}[cmap=UniGB-UTF16-H]
\setCJKmonofont{FandolFang-Regular.otf}[cmap=UniGB-UTF16-H]
\usepackage{amsmath,amssymb,mathtools,bm,amsthm}
\usepackage{booktabs,longtable,tabularx,array}
\usepackage{enumitem,fancyhdr,xcolor,microtype}
\usepackage[hidelinks,unicode]{hyperref}
\usepackage{bookmark}
\hypersetup{pdftitle={全向干扰源最优第二检测点的贝叶斯布站模型},pdfauthor={}}
\ctexset{section={format=\Large\bfseries,beforeskip=1.3ex plus .4ex,afterskip=.8ex},subsection={format=\large\bfseries,beforeskip=1ex plus .3ex,afterskip=.5ex}}
\setlength{\parindent}{2em}
\setlength{\parskip}{.22em}
\linespread{1.18}
\setlist[enumerate]{leftmargin=2.3em,itemsep=.3em,topsep=.4em}
\setlist[itemize]{leftmargin=2em,itemsep=.2em}
\pagestyle{fancy}
\fancyhf{}
\fancyhead[L]{\small 2026 年数学建模竞赛 B 题 · 问题二}
\fancyhead[R]{\small 模型与数值方法}
\fancyfoot[C]{\thepage}
\renewcommand{\headrulewidth}{.35pt}
\setlength{\emergencystretch}{2em}
\clubpenalty=10000
\widowpenalty=10000
\allowdisplaybreaks[1]
\numberwithin{equation}{section}
\theoremstyle{definition}
\newtheorem{proposition}{命题}[section]
\newtheorem{remark}{说明}[section]
\renewcommand{\proofname}{证明}
\newcommand{\R}{\mathbb R}
\newcommand{\E}{\mathbb E}
\newcommand{\Prob}{\mathbb P}
\newcommand{\ind}{\mathbf 1}
\newcommand{\dd}{\,\mathrm d}
\newcommand{\norm}[1]{\left\lVert #1\right\rVert}
\newcommand{\wrap}{\operatorname{wrap}}
\newcommand{\supp}{\operatorname{supp}}
\newcommand{\dist}{\operatorname{dist}}
\newcommand{\argmin}{\operatorname*{arg\,min}}
\newcommand{\argmax}{\operatorname*{arg\,max}}
\newcommand{\cl}{\operatorname{cl}}
\newcommand{\area}{\operatorname{area}}
\newcommand{\MEC}{\operatorname{MEC}}
\newcolumntype{Y}{>{\raggedright\arraybackslash}X}

\title{\sffamily 全向干扰源最优第二检测点的贝叶斯布站模型\\[.4em]\large 以第二次观测后最小包围圆的期望面积为目标}
\date{2026 年 9 月 11 日}
\begin{document}
\maketitle
\vspace{-2em}

\section{问题条件、研究任务与评价边界}
\subsection{题目提供的条件}
根据原题问题二及附录 1、附录 2\cite{problem2026b}，本研究对象为一个全向干扰源，已知一个检测点的坐标及该点测得的正常示向度，需要设计第二检测点的选择策略并给出候选区域。与本问直接相关的条件汇总于表~\ref{tab:facts}。源的总数、全区域多源搜索顺序和正式测试要求属于后续问题，不作为本问单源布站模型的额外输入。

\begin{table}[htbp]
\centering\small
\caption{题目条件及其在本模型中的作用}\label{tab:facts}
\begin{tabularx}{\textwidth}{p{.22\textwidth}Y}
\toprule
题目条件 & 模型中的含义 \\
\midrule
目标分布圆域半径为 1800 米 & 圆心为原点，正东为横轴正向、正北为纵轴正向；源位置限定在该圆域内。\\
全向源在全部方向辐射 & 接收与否只由源距和有效接收半径决定，不存在定向盲区。\\
误差位于 \([-1^\circ,1^\circ]\) & 每次正常示向度确定半角为 \(1^\circ\) 的方向约束；题目没有指定误差概率密度。\\
同一地点的误差在一段时间内固定 & 原地重复检测不能作为独立观测来压缩不确定性。\\
接收半径在 1000 至 1500 米之间 & 半径对同一个源固定，但数值未知；不能把上下界当成源距的上下界。\\
距离不超过 5 米时信号过强 & 无正常示向度，进入近距离强信号分支，可直接进行光学定位。\\
光学定位及清除有效距离为 20 米 & 若某个位置到所有可行源位置的距离均不超过 20 米，则可保证在该位置完成后续操作。\\
移动速度为 5 米每秒 & 用于定位效果接近时的移动距离或时间比较。\\
正常检测或确认无信号耗时 5 秒 & 第二次单源检测的固定时间项；切换频道为 1 秒，光学定位为 3 秒，清除为 2 秒。\\
\bottomrule
\end{tabularx}
\end{table}

题目未给出障碍物、不可达区域或第二点必须位于目标分布圆域内的约束。目标圆域限制的是源位置；除非另行给出移动限制，不能自动将其当作机器狗活动边界。本模型使用平面直线移动，并在候选域构造时说明有效搜索范围。

\subsection{从交叠面积到最小包围圆面积}
第一次示向度通常留下狭长的扇形区域，第二次测向的作用是缩小源的位置不确定性。若仅以交叠区域面积评价定位效果，一个很长但极窄的集合也可能获得较小损失。例如，\(80\,\mathrm m\times0.5\,\mathrm m\) 矩形的面积为 \(40\,\mathrm m^2\)，最小包围圆半径约为 \(40\,\mathrm m\)；\(10\,\mathrm m\times10\,\mathrm m\) 正方形的面积为 \(100\,\mathrm m^2\)，相应半径为 \(\sqrt{50}\,\mathrm m\)。前者面积更小，位置的最坏偏离却明显更大。

因此，本文以最小包围圆面积作为每次反馈后的损失。评价时刻固定为\textbf{第二次无线电检测刚结束、尚未执行光学定位时}。近距离强信号分支保留其半径不超过 5 米的可行区域，不在主模型中人为将该分支损失设为零。若将评价时刻改为完成强信号后的光学操作，只需将对应损失替换为零，但必须对所有比较策略使用相同口径。

本问以平均定位质量为主目标，以移动距离和达到 20 米保证定位条件的概率为辅助指标，不直接将时间与面积相加。

\section{假设、符号与定义域}
\subsection{全部建模假设及理由}
\begin{enumerate}[label=\textbf{A\arabic*.}]
\item \textbf{位置面积均匀先验。} 源位置 \(G\) 在目标圆域 \(D\) 内按面积均匀分布，故 \(p_G(g)=\ind_D(g)/|D|\)。该假设体现区域内无额外位置偏好，表示同面积的两块区域具有相同先验概率，不表示源距均匀。
\item \textbf{均匀角误差。} 在可获得正常示向度的条件下，给定源位置及接收半径，某个固定检测点在环境总体上的误差分布为 \(U[-\alpha,\alpha]\)，其中 \(\alpha=\pi/180\)。该假设将有界误差具体化为概率观测模型，并不由题目中的误差界自动推出。
\item \textbf{不同地点误差条件独立。} 给定 \(G,R\)，不同检测点的误差相互独立，且误差不影响接收阈值。该补充假设使第二次测向具有可计算的预测分布。把不同地点视为独立环境是基准近似，题目并未排除空间相关性；同一地点始终使用同一误差，不允许独立重抽。
\item \textbf{接收半径先验。} 基准模型取 \(R\sim U[R_-,R_+]\)，其中 \(R_-=1000\)、\(R_+=1500\)，并假设 \(R\) 与 \(G\) 独立。这是在仅知上下界时选取的透明、便于复现的参考先验，并非已知的真实生成机制。一般推导保留密度 \(h(r)\)，以便检验先验敏感性；同一源的 \(R\) 在整个观测过程中保持固定。
\end{enumerate}

所有概率是对未知环境参数的描述；在一次实际任务中，源位置、接收半径和各位置的环境误差均已有固定实现。推导不要求同一检测点的误差随时间随机变化。若发现位置与接收半径相关，应将独立先验替换为联合先验；若发现不同地点误差相关，应替换第二次条件误差分布。

\subsection{符号、输入与决策变量}
令 \(B(c,r)=\{x\in\R^2:\norm{x-c}\le r\}\)，并定义圆周角差 \(\wrap(t)=((t+\pi)\bmod 2\pi)-\pi\)。所有距离的单位为米、角度为弧度、面积为平方米。

\begin{longtable}{p{.24\textwidth}p{.36\textwidth}p{.31\textwidth}}
\caption{主要符号及定义域}\label{tab:symbols}\\
\toprule 符号 & 含义 & 定义域或取值 \\
\midrule\endfirsthead
\toprule 符号 & 含义 & 定义域或取值 \\
\midrule\endhead
\bottomrule\endfoot
\(D\) & 目标位置先验区域 & \(B(0,1800)\) \\
\(S_1,\theta_1\) & 第一点及已知正常示向度 & \(S_1\in\R^2\)，\(\theta_1\in[0,2\pi)\) \\
\(G,R\) & 未知源位置、固定接收半径 & \(G\in D\)，\(R\in[1000,1500]\) \\
\(q=(q_x,q_y)\) & 待优化的第二检测点 & \(q\in\mathcal C\subset\R^2\) \\
\(r_0,r_c,\alpha\) & 强信号阈值、光学阈值、半角 & \(5,20,\pi/180\) \\
\(r_1(g),r_2(g;q)\) & 候选源至第一、第二点的距离 & \(\norm{g-S_1},\norm{g-q}\) \\
\(\phi_s(g)\) & 从检测点 \(s\) 指向 \(g\) 的方位 & \(\arg(g-s)\bmod2\pi\)，\(g\ne s\) \\
\(\mathcal I_1\) & 第一次反馈信息 & 收到正常示向度 \(\theta_1\) \\
\(Z\) & 第二次反馈 & \(\{H,N\}\sqcup[0,2\pi)\) \\
\(H,N\) & 近距离强信号、无信号标签 & 两个离散反馈值 \\
\(p_1,p_2\) & 第一次、第二次后的位置密度 & 对平面面积测度的概率密度 \\
\(K_1,K_z(q)\) & 第一次、反馈 \(z\) 后的位置支持集 & 闭、有界；零概率分支不计损失 \\
\(c(K),\rho(K)\) & 最小包围圆圆心、半径 & \(c(K)\in\R^2\)，\(\rho(K)\ge0\) \\
\(J(q)\) & 第二次后包围圆期望面积 & \(\mathrm m^2\) \\
\(P_{20}(q)\) & 达到 20 米保证定位条件的概率 & \([0,1]\) \\
\(b_0,\eta,\varepsilon_{\rm abs}\) & 布站间距、相对和绝对面积容差 & \(b_0>0\)，\(\eta\ge0\)，\(\varepsilon_{\rm abs}\ge0\) \\
\end{longtable}

独立自变量为 \(q_x,q_y\)，或与其等价的基线长度和偏转角。\(G,R,Z\) 是要积分掉的不确定量，不能在布站时使用其真实值。反馈后的圆心是依赖实际反馈的后续变量，不能提前对所有反馈强制使用同一圆心。

\section{第一次正常示向度的贝叶斯更新}
\subsection{正常观测的似然}
记 \(a_1(g)=\ind\{|\wrap(\theta_1-\phi_{S_1}(g))|\le\alpha\}\)。给定 \(G=g,R=r\)，第一次返回正常示向度 \(\theta_1\) 的似然密度为 \(\ell_1(g,r)=\ind\{r_0<r_1(g)\le r\}a_1(g)/(2\alpha)\)。该式同时包括正常接收条件与测角误差条件。

为避免将零测度边界误当成独立的可能位置，定义正密度集合 \(\Omega_1=\{g\in D:r_0<r_1(g)<R_+,\ a_1(g)=1\}\)，并在几何计算中取其概率支持集 \(K_1=\supp p_1\)。在非退化输入下，它等于 \(\Omega_1\) 中正面积部分的闭包。连续概率计算中边界不改变积分，但最小包围圆必须覆盖支持集的极限边界。

\subsection{联合后验及接收概率函数}
定义半径先验的尾概率 \(w(d)=\int_{R_-}^{R_+}\ind\{r\ge d\}h(r)\dd r\)。在基准均匀先验下，该函数为
\begin{equation}\label{eq:w}
w(d)=
\begin{cases}
1,&0\le d\le R_-,\\
(R_+-d)/(R_+-R_-),&R_-<d<R_+,\\
0,&d\ge R_+.
\end{cases}
\end{equation}
由贝叶斯公式，先验中的常数 \(1/|D|\) 及正常角度似然中的常数 \(1/(2\alpha)\) 同时出现在分子和分母，因而消去。设 \(Z_1=\int_{\Omega_1}w(r_1(g))\dd g\)，则联合后验与位置边缘后验分别满足
\begin{equation}\label{eq:posterior1}
\begin{aligned}
p(g,r\mid\mathcal I_1)
&=\frac{\ind_{\Omega_1}(g)\,\ind\{r\ge r_1(g)\}\,h(r)}{Z_1},\\
p_1(g)&=\frac{\ind_{\Omega_1}(g)\,w(r_1(g))}{Z_1}.
\end{aligned}
\end{equation}
要求输入满足 \(Z_1>0\)，否则第一次观测与模型不相容，不能继续归一化。进一步，给定 \(G=g\) 及第一次反馈后，\(R\) 的条件密度为 \(h(r)\ind\{r\ge r_1(g)\}/w(r_1(g))\)。均匀先验下，它是区间 \([\max\{R_-,r_1(g)\},R_+]\) 上的均匀分布。

式~\eqref{eq:posterior1} 表明，第一次正常接收已经使 \(G\) 与 \(R\) 后验相关，尽管二者的先验独立。源距大于 1000 米时，成功接收需要更大的半径，因此其位置权重降低。若只利用示向度、不利用成功接收信息，均匀位置先验和恒定角度似然确实给出截断扇形内的面积均匀后验；此处出现的非均匀性来自接收事件，而不是从角度误差中额外推断出了距离。

例如，完整环形扇形上的面积均匀分布在极坐标中含有面积元 \(r\dd r\dd\varphi\)，径向密度与 \(r\) 成正比。因此，即使暂时忽略接收信息，也不能在距离和角度上分别均匀取样，再将样本视为平面面积均匀。

\section{第二次反馈、预测分布与位置支持集}
\subsection{完整反馈空间}
给定第二点 \(q\)，第二次反馈 \(Z\) 是一个离散与连续混合的观测。若 \(r_2(g;q)\le r_0\)，返回 \(H\)；若 \(r_2(g;q)>r\)，返回 \(N\)；若 \(r_0<r_2(g;q)\le r\)，返回正常示向度 \(\Theta=(\phi_q(g)+E_2)\bmod2\pi\)。由于 \(R_-\gg r_0\)，强信号与无信号分支不会重叠。连续角度使用勒贝格测度，两个离散标签使用计数测度。

\subsection{积分掉同一个固定接收半径}
以下用 \(r_2=r_2(g;q)\)、\(r_1=r_1(g)\) 简化符号，并令 \(m(g;q)=\max\{r_1,r_2\}\)。定义未归一化位置权重
\begin{equation}\label{eq:bweights}
\begin{aligned}
b_H(g;q)&=\ind_{\Omega_1}(g)\ind\{r_2\le r_0\}\,w(r_1),\\
b_N(g;q)&=\ind_{\Omega_1}(g)\,[w(r_1)-w(m(g;q))],\\
b_\theta(g;q)&=\frac{\ind_{\Omega_1}(g)\ind\{r_2>r_0\}}{2\alpha}
\ind\{|\wrap(\theta-\phi_q(g))|\le\alpha\}\,w(m(g;q)).
\end{aligned}
\end{equation}
第一式来自第一次成功接收条件与第二次近距离条件。第二式对应 \(r_1\le R<r_2\)，当 \(r_2\le r_1\) 时自然为零。第三式对应两次都能正常接收，必须由\textbf{同一个}半径满足 \(R\ge\max\{r_1,r_2\}\)，再乘第二次的均匀角度似然。这里不能使用 \(w(r_1)w(r_2)\)，也不能在第二次观测前重新从原始半径先验独立抽取 \(R\)。

对 \(z=H,N\) 及正常角度 \(z=\theta\)，令 \(M_z(q)=\int_D b_z(g;q)\dd g\)。相应预测概率、预测密度及位置后验为
\begin{equation}\label{eq:prediction}
\begin{gathered}
P_H(q)=\frac{M_H(q)}{Z_1},\qquad
P_N(q)=\frac{M_N(q)}{Z_1},\qquad
f_\Theta(\theta;q)=\frac{M_\theta(q)}{Z_1},\\
p_2(g\mid z,q,\mathcal I_1)=\frac{b_z(g;q)}{M_z(q)}\quad\text{当 }M_z(q)>0.
\end{gathered}
\end{equation}
其中 \(f_\Theta\) 是正常反馈的未条件化密度，其积分为正常反馈发生概率，而不是恒为 1。若先除以正常反馈发生概率再计算损失，必须在总目标中乘回该概率，否则会隐去第二点失去信号的风险。

\begin{proposition}[反馈概率守恒]\label{prop:normalization}
在本模型下，\(P_H(q)+P_N(q)+\int_0^{2\pi}f_\Theta(\theta;q)\dd\theta=1\)。
\end{proposition}
\begin{proof}
对于固定的 \(g\in\Omega_1\)，正常角度似然在圆周上积分为 1。若 \(r_2>r_0\)，式~\eqref{eq:bweights} 中正常反馈积分为 \(w(m)\)，无信号权重为 \(w(r_1)-w(m)\)，二者之和为 \(w(r_1)\)。若 \(r_2\le r_0<r_1\)，则 \(m=r_1\)，无信号和正常反馈权重均为零，强信号权重为 \(w(r_1)\)。两种情况下，总权重都等于第一次后验的未归一化权重。对 \(g\) 积分再除以 \(Z_1\)，即得结论。
\end{proof}

\subsection{支持集的几何表达}
对 \(M_z(q)>0\) 的反馈，定义 \(K_z(q)=\supp p_2(\cdot\mid z,q,\mathcal I_1)\)。在基准半径先验下，\(h\) 在 \((R_-,R_+)\) 内为正，其支持集可由下列正权重区域取闭包得到
\begin{equation}\label{eq:supports}
\begin{aligned}
K_H(q)&=\cl\{g\in\Omega_1:r_2(g;q)<r_0\},\\
K_N(q)&=\cl\{g\in\Omega_1:r_2(g;q)>\max\{R_-,r_1(g)\}\},\\
K_\theta(q)&=\cl\{g\in\Omega_1:r_0<r_2(g;q)<R_+,\
|\wrap(\theta-\phi_q(g))|\le\alpha\}.
\end{aligned}
\end{equation}
式中取闭包时只保留正面积可行部分的极限点；相切形成的孤立零面积分支不会产生独立的正概率后验。对于零概率离散反馈或零密度角度节点，可将其积分贡献设为零，而不对空集计算包围圆。

正常反馈区域是两组角度约束、目标圆域及接收距离约束的交集，并排除近距离圆盘内部；它未必是图示中的纯四边形。无信号区域也不应直接认定为第一次区域不变，它包含第一次接收与第二次未接收之间的距离比较信息。由于 \(r_2^2-r_1^2=2g\cdot(S_1-q)+\norm{q}^2-\norm{S_1}^2\)，条件 \(r_2>r_1\) 实际上是一条直线限定的半平面。这一性质使无信号区域仍可用直线和圆弧边界处理。

式~\eqref{eq:supports} 对其他在整个半径区间内正密度的先验仍成立，但概率权重会改变。若半径已知、具有点质量或先验支持区间缩小，应从式~\eqref{eq:bweights} 的正权重集合重新构造支持集，不能机械沿用上述严格不等式。

\section{最小包围圆损失及其定位含义}
\subsection{反馈后的内层优化}
对非空紧集 \(K\subset D\)，定义其最小包围圆半径与圆心为
\begin{equation}\label{eq:mec}
\rho(K)=\min_{c\in\R^2}\ \max_{g\in K}\norm{g-c},
\qquad
c(K)\in\argmin_{c\in\R^2}\ \max_{g\in K}\norm{g-c}.
\end{equation}
该最小值存在，且 \(\rho(K)\le1800\)，因为以原点为圆心、半径 1800 米的圆已经覆盖 \(D\)。圆盘是凸集，因此覆盖 \(K\) 等价于覆盖其凸包；但计算时不能仅用有限样本的凸包代替整个连续可行区域。

令 \(\rho_z(q)=\rho(K_z(q))\)、\(c_z(q)=c(K_z(q))\)，损失定义为 \(L_z(q)=\pi\rho_z(q)^2\)。该损失的含义是，获得反馈后选择最有利的估计位置，使最坏位置误差最小，再以此误差半径所对应的圆面积评价定位质量。圆心允许随反馈变化，因此数学顺序是先对每个反馈求内层最优，再对反馈取期望，而不是先固定圆心再优化。

\begin{proposition}[光学定位保证]
若 \(\rho_z(q)\le r_c=20\)，机器人到达 \(c_z(q)\) 后，可对支持集内的任意真实源位置保证完成光学定位与清除。
\end{proposition}
\begin{proof}
由式~\eqref{eq:mec}，对每个 \(g\in K_z(q)\) 都有 \(\norm{g-c_z(q)}\le\rho_z(q)\le20\)。题目规定光学操作只依赖该距离，因此结论成立。
\end{proof}

对固定区域，最小化圆面积与最小化半径等价；对随机反馈，最小化 \(\E[\rho^2]\) 与最小化 \(\E[\rho]\) 一般不等价。本文选择前者，以较重地惩罚少量较差反馈导致的大定位半径。最小包围圆覆盖所有仍然可能的位置，低后验密度的边缘位置也不能被忽略，这是获得支持集层面保证所需的保守性。

\subsection{期望面积目标}
将式~\eqref{eq:prediction} 的预测分布与式~\eqref{eq:mec} 的内层优化结合，得到第二点的目标函数
\begin{equation}\label{eq:objective}
\boxed{\displaystyle
J(q)=\frac{\pi}{Z_1}\left[
M_H(q)\rho_H(q)^2+M_N(q)\rho_N(q)^2
+\int_0^{2\pi}M_\theta(q)\rho_\theta(q)^2\dd\theta
\right].}
\end{equation}
该函数依赖 \(q_x,q_y\)，并依赖已知输入 \(S_1,\theta_1\) 及假设参数 \(h,\alpha\)。\(J\) 的量纲为平方米，且 \(0\le J(q)\le\pi\rho(K_1)^2\le\pi1800^2\)，因为每种正概率反馈的支持集均为 \(K_1\) 的子集。强信号分支满足 \(\rho_H(q)\le5\)，但按照本文的评价时刻，仍使用其实际包围圆面积。

与蒙特卡洛形式对应，式~\eqref{eq:objective} 也可写为 \(J(q)=\E_{G,R\mid\mathcal I_1}\E_{E_2}[L_{Z(q;G,R,E_2)}(q)]\)。这个形式揭示平均所针对的全部不确定性，式~\eqref{eq:objective} 则已解析消去了第二次误差和半径的一部分计算，通常更适合稳定积分。

\section{约束、候选区域与完整优化问题}
\subsection{全反馈模型的有用搜索域}
如果 \(\dist(q,K_1)>R_+\)，任何可行源均不可能在第二点被接收到，此时 \(N\) 为确定反馈且没有新增信息，\(J(q)=\pi\rho(K_1)^2\)。因此，在寻找有信息的第二点时，只需搜索
\begin{equation}\label{eq:candidate}
\mathcal C=\{q\in\R^2:\norm{q-S_1}\ge b_0,\ \dist(q,K_1)\le R_+\},
\qquad
\boxed{q^\star\in\argmin_{q\in\mathcal C}J(q).}
\end{equation}
对所选 \(b_0\) 要先检查 \(\mathcal C\ne\varnothing\)。由 \(K_1\subset B(0,1800)\)，搜索域包含于 \(B(0,3300)\)。边界上只有相切、零接收概率的点可作为无信息基准保留，不影响有信息点的筛选。若在强移动约束下有用搜索域为空，应报告该约束下无法获得第二次有信息检测，而不是强行输出坐标。

为了便于解释布站方向，令 \(u_1=(\cos\theta_1,\sin\theta_1)\)、\(v_1=(-\sin\theta_1,\cos\theta_1)\)，可将 \(q\) 写为 \(S_1+b(\cos\beta\,u_1+\sin\beta\,v_1)\)。等价自变量为 \(b\ge b_0\)、\(\beta\in[-\pi,\pi)\)，还需满足式~\eqref{eq:candidate} 的距离约束。目标圆域仍以原点为中心，坐标转换后不能忽略其截断效应。

第二点与第一点完全相同时，同一地点误差不变，反馈是已知信息，必须赋值 \(J(S_1)=\pi\rho(K_1)^2\)，不能使用独立误差公式制造额外信息。本模型通过 \(b_0>0\) 排除此情况。布站间距参数用于形成可复现的约束问题，不构成对误差空间相关长度的估计。

\subsection{保证继续接收的可选约束}
第一次收到信号意味着，对真实源位置 \(g\)，有 \(R\ge\max\{R_-,r_1(g)\}\)。据此定义
\begin{equation}\label{eq:csafe}
\mathcal C_{\rm safe}=\left\{q\in\mathcal C:
\norm{q-g}\le\max\{R_-,r_1(g)\},\quad\forall g\in K_1\right\}.
\end{equation}
在该区域内，第二次不会因距离超出有效接收半径而返回无信号，但可能进入有利的强信号分支。该条件比 \(\max_{g\in K_1}\norm{q-g}\le1000\) 更宽松，因为已经利用了第一次接收给出的半径下界。

式~\eqref{eq:candidate} 是允许全部反馈的主模型；将其可行域换成 \(\mathcal C_{\rm safe}\) 得到保证接收的受限模型。两者应分别求解并比较，后者不能被称为不带该约束的全局最优。若 \(\mathcal C_{\rm safe}\) 为空，主模型仍可以使用。

\subsection{近优候选区及辅助排序}
设 \(J_{\min}\) 为主模型的最小目标值，定义近优候选区域为 \(\mathcal C_\eta=\{q\in\mathcal C:J(q)\le J_{\min}+\varepsilon_{\rm abs}+\eta J_{\min}\}\)。加入绝对容差是为了在最小值接近零时仍有清晰的数值标准。\(\eta\) 是允许的相对性能损失，不是由题目给定的概率。

第二次后达到保证光学定位条件的概率为
\begin{equation}\label{eq:p20}
\begin{aligned}
P_{20}(q)={}&P_H(q)\ind\{\rho_H(q)\le20\}
+P_N(q)\ind\{\rho_N(q)\le20\}\\
&+\int_0^{2\pi}f_\Theta(\theta;q)\ind\{\rho_\theta(q)\le20\}\dd\theta.
\end{aligned}
\end{equation}
它表示能够选择一个点保证清除的反馈概率，不等于所有可能清除策略的实际成功概率。即使某个支持集的包围圆半径大于 20 米，也可能存在某个位置覆盖较高的后验概率质量。

在主指标相近时，可以先比较 \(P_{20}\)，再选择移动距离较小的点。若频道保持不变，则移动加第二次检测耗时为 \(\norm{q-S_1}/5+5\) 秒；如果必须切换频道，增加一个与 \(q\) 无关的 1 秒常数。该时间尚未包含从第二点到后验圆心的后续移动，不能标为完整清除总时间。

\section{数值求解方法}
\subsection{计算结构与输入输出}
输入为 \(S_1,\theta_1\)、接收半径先验 \(h\)、布站参数 \(b_0\)、目标容差及几何和积分精度。输出包括近似最优第二点、近优候选区域、期望包围圆面积、\(P_H,P_N,P_{20}\)、移动距离及数值误差记录。对具体反馈还应能输出位置支持集和包围圆圆心、半径，以便随后执行。

计算分为四层：空间积分得到预测权重，连续几何优化得到各反馈的包围圆，反馈积分得到 \(J(q)\)，二维搜索得到候选点。\(G\) 的加权样本可以用于第一层和独立模拟，但\textbf{不能把有限随机样本的最小包围圆当作连续支持集的最小包围圆}。

\subsection{第一次后验与空间积分}
用以 \(S_1\) 为原点的极坐标表示候选源位置，令 \(g=S_1+r(\cos\varphi,\sin\varphi)\)。第一扇形限定 \(\varphi\in[\theta_1-\alpha,\theta_1+\alpha]\)，距离限定 \(r\in(r_0,R_+)\)。任意可积函数 \(F\) 的空间积分可写为
\begin{equation}\label{eq:polar}
\int_{\Omega_1}F(g)\dd g
=\int_{\theta_1-\alpha}^{\theta_1+\alpha}
\int_{r_0}^{R_+}
F\!\left(S_1+r(\cos\varphi,\sin\varphi)\right)
\ind_D\!\left(S_1+r(\cos\varphi,\sin\varphi)\right)
\,r\dd r\dd\varphi.
\end{equation}
内层的圆域截断可通过 \(\norm{S_1+ru(\varphi)}^2\le1800^2\) 的二次不等式精确求得径向区间。计算 \(M_z\) 时，还要沿射线求解第二点距离圆和角度边界的交点，把积分区间在这些边界处分段。这样可以避免在不连续指示函数上盲目使用高阶求积。

推荐用自适应区间求积或自适应二维三角剖分，保持面积权重非负。粗搜索可预先生成等面积网格与权重，精搜索再使用自适应积分。若使用普通极坐标网格，每个单元的面积必须按 \((r_{i+1}^2-r_i^2)(\varphi_{j+1}-\varphi_j)/2\) 计算；边界单元还需裁剪或细分。

\subsection{连续可行区域的边界构造}
式~\eqref{eq:supports} 涉及圆盘、圆盘外部和半平面的交集。小于 \(\pi\) 的方向扇形可以用两条射线对应的半平面表示；无信号条件 \(r_2>r_1\) 已转化为线性不等式。因此，边界由有限条线段和圆弧组成，可能具有多个连通分支或孔洞。

边界构造时，先计算所有相关直线、圆之间的交点，再以交点划分线段和圆弧，在各开边界片段上检验其邻域是否与正面积可行部分相邻，保留支持集的真实边界。对无交点的整圆或完整可行边界分支单独判定；角度跨越 \(0\) 时用向量半平面处理，避免角度排序错误。仅满足等号但不邻接任何正密度可行区域的孤立分量应排除。

若改用多边形近似圆弧，需要明确内接或外接的方向及最大几何误差。内接近似一般低估支持范围，不能直接提供清除保证；用于保证定位的输出必须采用经认证的外包络或随后进行连续边界覆盖检验。

\subsection{最小包围圆的约束生成与上下界}
对固定的非空 \(K=K_z(q)\)，式~\eqref{eq:mec} 是关于 \(c\in\R^2\)、\(\rho\ge0\) 的半无限凸约束问题，约束为 \(\norm{g-c}\le\rho\) 对所有 \(g\in K\) 成立。采用最远点约束生成可在每次迭代获得可核查的半径上下界。

首先从 \(K\) 选取有限非空点集 \(V\subset K\)，计算其最小包围圆 \((c_V,\rho_V)\)。有限点问题可用二阶锥规划求解；也可使用由一个点、两个点的直径圆或三个不共线点的外接圆确定的最小圆算法，并检查其覆盖全部有限点。由 \(V\subset K\)，有下界 \(\rho_V\le\rho(K)\)。随后在完整集合上计算 \(g_{\rm far}\in\argmax_{g\in K}\norm{g-c_V}\)，记 \(d_V=\norm{g_{\rm far}-c_V}\)。同一圆心半径 \(d_V\) 覆盖 \(K\)，所以
\begin{equation}\label{eq:mecbounds}
\rho_V\le\rho(K)\le d_V.
\end{equation}
若 \(d_V-\rho_V\le\varepsilon_\rho\)，即可停止；否则将 \(g_{\rm far}\) 加入 \(V\) 后重新求有限点最小圆。数值实现中，应把有限凸问题的求解误差和最远点计算误差计入上下界，不能把浮点近似值当成未经修正的严格界。

最远点搜索可直接利用连续边界。在线段上，距离平方是凸函数，最大值在端点取得。在圆心为 \(o\)、半径为 \(a\) 的圆弧上，点写为 \(o+a(\cos t,\sin t)\)，其到 \(c_V\) 的距离平方为 \(\norm{o-c_V}^2+a^2+2a(o-c_V)\cdot(\cos t,\sin t)\)。因此只需检查弧段端点，以及方向 \(\arg(o-c_V)\) 是否位于该弧段内；若 \(o=c_V\)，整段距离相同。遍历所有真实边界片段即可得到全局最远距离，无需用密集随机点代替连续覆盖检验。

上下界对应面积区间 \([\pi\rho_V^2,\pi d_V^2]\)。若停止条件为 \(d_V-\rho_V\le\varepsilon_\rho\)，并使用 \(\rho(K)\le1800\)，面积差不超过 \(\pi(3600\varepsilon_\rho+\varepsilon_\rho^2)\)。这是统一的保守界，实际应记录每个反馈的 \(\pi(d_V^2-\rho_V^2)\) 并按反馈概率加权。输出保证覆盖圆时，可使用 \((c_V,d_V)\)；如其半径不超过 20 米，则清除保证成立。若下界已经超过 20 米，该反馈不满足保证条件；上下界跨越 20 米时需继续加密。

\subsection{对第二次示向度积分}
先计算两个离散反馈分支的 \(M_H,\rho_H,M_N,\rho_N\)，再对 \(\theta\in[0,2\pi)\) 计算 \(M_\theta\rho_\theta^2\) 的积分。只有从 \(q\) 观察 \(\Omega_1\cap\{r_0<r_2<R_+\}\) 的方位范围向两侧扩张 \(\alpha\) 后，才可能出现正常示向度。实际实现应按圆周上的一个或多个区间处理；当候选点在可行区域附近时，角度范围可能很宽，不能始终假设只有一个狭窄区间。

推荐在可行角度区间内使用自适应求积，并在支持集拓扑变化或包围圆支撑点切换附近细分。粗细两级积分差用于误差估计，但对不连续指标 \(\ind\{\rho_\theta\le20\}\)，还需定位阈值跨越点或继续细分，不应沿用对平滑函数的误差判断。

若为节省计算忽略一组低概率角度，其总预测概率为 \(\delta\)，则被忽略的目标贡献最多为 \(\pi1800^2\delta\)。这给出可说明的截断误差界；不能仅因为某个角度节点概率小就把其可能很大的包围圆损失任意设为零。空间积分、角度积分和包围圆误差应分别报告。

\subsection{第二点的二维搜索}
目标 \(J(q)\) 涉及集合形状变化和最小包围圆支撑点变化，通常不应假定为凸函数或处处可微。推荐在式~\eqref{eq:candidate} 的有界域内进行多层搜索：先覆盖完整候选域，保留多个相互分离的较优区域，再细化网格并进行无导数局部搜索，最后以更严格的积分和几何精度重新评价入选点。

在候选域边界、保证接收区域边界以及首次示向度两侧应保留搜索点，不应只在一个主观选定的侧向扇区内优化。近优区域可能有多个连通分支，尤其在对称输入下常出现成对解。仅搜索一侧必须有严格的对称性依据。

有限网格上全部点的最优值可以被直接验证，但它只是该网格的最优值。网格加密和多起点一致性提供数值稳定证据，不能自动证明连续域的全局最优。若需要全局证书，必须进一步给出每个布站网格单元的目标下界并使用分支定界；本文没有预设未经证明的全局 Lipschitz 常数，因此不把常规网格收敛称为全局最优性证明。

\begin{table}[htbp]
\centering\small
\caption{建议的分层数值控制与报告内容}\label{tab:numeric}
\begin{tabularx}{\textwidth}{p{.20\textwidth}Y Y}
\toprule 层次 & 控制方式 & 必须记录的量 \\
\midrule
源位置积分 & 边界分段、自适应求积、面积权重 & \(Z_1\)、各 \(M_z\) 及积分误差估计 \\
连续包围圆 & 有限点圆加连续最远点约束生成 & 半径上下界、覆盖残差、迭代次数 \\
反馈积分 & 有效角度区间、自适应细分 & 概率守恒残差、\(J\) 及 \(P_{20}\) 精度 \\
第二点搜索 & 全域粗网格、多区域细化、多起点 & 网格尺度、候选点坐标、目标稳定性 \\
假设敏感性 & 半径先验与布站间距变化 & 最优区域位移、目标和保证概率变化 \\
\bottomrule
\end{tabularx}
\end{table}

\subsection{完整算法流程}
\begin{enumerate}[label=步骤\arabic*：]
\item 读取输入，统一角度单位并构造第一扇形、目标圆域及距离条件；计算 \(Z_1\)，若其非正或与零无法数值区分，则报告输入或模型不相容。
\item 构造 \(K_1\)、主候选域 \(\mathcal C\) 及可选的 \(\mathcal C_{\rm safe}\)，确定 \(b_0\)、搜索网格与各层误差容差。
\item 对每个候选点 \(q\)，由式~\eqref{eq:bweights} 计算强信号和无信号的预测质量，并构造其支持集；对正概率分支，用连续最远点约束生成计算包围圆上下界。
\item 确定正常示向度的有效角度区间；在自适应角度节点上计算 \(M_\theta\)、\(K_\theta\) 和包围圆上下界，完成式~\eqref{eq:objective} 与式~\eqref{eq:p20} 的积分。
\item 检查反馈总概率、支持集包含关系与包围圆覆盖残差。出现不一致时提高积分或几何精度，未通过检查的候选点不进入最终排名。
\item 保留多个较优搜索区域并加密，对最终候选重新高精度评价；按近优阈值形成 \(\mathcal C_\eta\)，再比较 \(P_{20}\) 和移动距离。
\item 用独立的联合后验模拟验证最终候选的目标值，开展半径先验及布站间距敏感性分析，输出候选区域与数值精度说明。
\end{enumerate}

为了避免积分噪声改变候选排名，两个点的目标差若小于其数值误差之和，应视为无法可靠区分并继续加密。若能得到候选点目标区间 \([J_i^-,J_i^+]\)，可用 \(\min_iJ_i^+\) 作为网格最优值的上界；只有当某点下界已显著高于该上界及近优容差时，才可靠地排除该点。

\section{独立验证、敏感性与结果解释}
\subsection{不同表示之间的独立验证}
主计算通过式~\eqref{eq:bweights} 对半径和反馈概率进行积分，独立验证则直接生成完整观测。先从位置后验 \(p_1\) 抽取 \(G^{(i)}\)，再从 \(R\mid G^{(i)},\mathcal I_1\) 抽取同一个源的半径 \(R^{(i)}\)，最后在正常接收时抽取第二次误差并生成 \(Z^{(i)}\)。均匀半径先验下，第二步的区间必须是 \([\max\{1000,r_1(G^{(i)})\},1500]\)，不能返回原始 \([1000,1500]\) 区间。

对每个模拟反馈重新构造完整支持集，计算 \(L_i=\pi\rho(K_{Z^{(i)}}(q))^2\)，得到样本均值 \(\widehat J_{\rm MC}=n^{-1}\sum_iL_i\) 及标准误差 \(s_L/\sqrt n\)。比较它与确定性积分值，同时对比三类反馈的模拟频率和理论预测概率。模拟均值误差条用于检验统计一致性，不取代连续支持集的几何覆盖认证。

包围圆算法应另行用点、线段、矩形、等边三角形和圆盘等可手算对象检查。二维有限点凸优化与连续圆弧最远点检验采用不同表示，可以发现遗漏弧段、错误内接近似或支持点不足等问题；只把同一公式换一个程序重算不构成充分的独立验证。

\subsection{必须通过的性质检查}
概率检查包括 \(\int p_1=1\)、命题~\ref{prop:normalization} 及所有预测质量非负。集合检查包括 \(K_z(q)\subseteq K_1\subseteq D\)；损失检查包括 \(J(q)\le\pi\rho(K_1)^2\)、强信号半径不超过 5 米，以及完全无信息点的目标等于第一次包围圆面积。几何检查包括最终覆盖残差、圆心到最远边界的距离和上下界间隙。

对称性检查要同时变换目标域与首次观测。由于 \(D\) 是以原点为中心的圆域，将 \(S_1,q\) 同时绕原点旋转、并将 \(\theta_1\) 增加同一角度，目标值应保持不变。只有整个输入关于首次示向度轴对称时，轴两侧候选点才必然等价；\(S_1\) 偏离原点时不能直接假设这种镜像对称。

在 \(\alpha\) 逐渐减小、两次视线非退化且均正常接收的情形下，正常反馈支持集应收缩，定位损失趋于零；但若仍存在正概率无信号分支，不能据此宣称总目标也趋于零。第二点靠近第一点时的表现还依赖空间独立性假设，不能把数学上的独立均匀噪声近点极限当成真实环境必然规律。

\subsection{敏感性分析与适用限制}
对接收半径先验，至少比较均匀先验与在整个区间内正密度、分别偏向小半径和大半径的先验。比较时保留同一输入和数值精度，报告近优区域位置、\(J\)、\(P_N\) 与 \(P_{20}\) 的变化。如果引入已知半径或收缩支持区间的先验，须同步重建支持集。

不同地点误差的独立性是另一项关键风险。若给定相关误差模型，第二次误差密度应条件于第一次信息，此时式~\eqref{eq:bweights} 中的恒定均匀角似然需替换，支持集也可能改变。通过改变 \(b_0\) 只能检查布站约束敏感性，不能验证误差独立性本身。

主目标对反馈取平均，因此无法保证每一次的定位半径都小于 20 米。若任务另有可靠性要求，可以增加 \(P_{20}(q)\ge p_0\) 的机会约束，其中 \(p_0\) 必须明确给定；若该约束不可行，应报告不可行，而不能降低阈值后仍声称满足原要求。对于低概率边缘位置过度主导包围圆的情形，可另建高后验概率区域的包围圆模型，但那将放弃对全部可行位置的覆盖保证，不能与本文目标混用。

\section{模型结论与应提交的结果}
本模型用第一次正常接收及示向度构造联合后验，将第二点视为在观测发生前选择的二维决策变量，并对第二次全部可能反馈分别求最小包围圆后再取期望。最终主问题由式~\eqref{eq:posterior1}、\eqref{eq:bweights}、\eqref{eq:mec}、\eqref{eq:objective} 和约束~\eqref{eq:candidate} 完整定义。保证接收模型是增加约束~\eqref{eq:csafe} 后的可选比较方案，不能代替主模型而不说明约束变化。

最优第二点通常需要在横向交会、接近可能源位置和保持接收能力之间权衡。第一次方向约束留下的主要不确定性往往沿视线方向分布，适当侧移能够增加对距离的约束；但移动也会改变源距和接收概率，因此不能仅依据某个代表位置的直角交会给出统一最优坐标。边界截断、接收半径先验与首次检测位置都会改变最优候选区域。

数值结果应至少给出首次观测输入、假设参数、第二点坐标、主目标与数值误差、三类反馈概率、\(P_{20}\)、移动距离和近优区域，并说明是否施加保证接收约束。附带网格、积分和包围圆收敛记录，才能使近似最优结论可复核。由于本稿没有指定具体 \(S_1,\theta_1\)，本文给出的是可直接据此实现的完整模型和数值求解规程，不宣称已经得到某一实际实例的最优坐标或连续域全局最优证书。

\begin{thebibliography}{9}
\bibitem{problem2026b} 全国大学生数学建模竞赛组委会. 2026 年高教社杯全国大学生数学建模竞赛 B 题：无线电干扰源的快速自动定位与清除[Z]. 2026. 本稿依据本地题目 PDF 的问题二及附录 1、附录 2，未采用参考论文中的数值答案。
\end{thebibliography}
\end{document}
